\chapter{Mucormycosis}

\medskip
\noindent\textit{\textbf{Under review.} This chapter is an AI-assisted educational draft; it carries no source citations and is not a substitute for professional medical advice.}

\section*{Definition and Overview}
Mucormycosis (previously zygomycosis) is a rare but aggressive, rapidly progressive fungal infection caused by moulds of the order Mucorales (including \textit{Rhizopus}, \textit{Mucor}, \textit{Lichtheimia}, and \textit{ Cunninghamella} species). These fungi are ubiquitous saprophytes found in soil, decaying organic matter, and bread/vegetable mould. Infection is acquired primarily through inhalation of sporangiospores, but may also follow direct cutaneous inoculation following trauma or burns. The disease predominantly affects immunocompromised hosts, including those with uncontrolled diabetes, haematological malignancy, and solid organ or haematopoietic stem cell transplants. Clinically it manifests in several forms, most commonly rhino-orbital-cerebral, pulmonary, cutaneous, gastrointestinal, and disseminated disease, each carrying a high mortality if not treated urgently.

\section*{Epidemiology}
Mucormycosis is a rare disease but has been increasing in incidence as the population of immunosuppressed patients grows. The global incidence is estimated at roughly 0.2--1.7 cases per million persons per year, though it is markedly higher among specific risk groups. In India, an outbreak of rhino-orbital-cerebral mucormycosis ("black fungus") occurred during the COVID-19 pandemic in 2021, driven by uncontrolled diabetes and misuse of corticosteroids. In Australia, mucormycosis is uncommon; most reported cases occur in diabetics with ketoacidosis, patients with haematological malignancy receiving intensive chemotherapy, transplant recipients, and patients on long-term deferoxamine or broad-spectrum antifungals. Mortality ranges from 20--30\% for localised cutaneous disease to 50--90\% for pulmonary and rhino-orbital-cerebral forms.

\section*{Aetiology and Pathophysiology}
\begin{itemize}
    \item \textbf{Pathogens:} Mucorales moulds, most commonly \textit{Rhizopus oryzae}; others include \textit{Rhizopus microsporus}, \textit{Mucor}, \textit{Lichtheimia corymbifera}, \textit{Cunninghamella}, and \textit{Saksenaea}
    \item \textbf{Predisposing factors:} diabetic ketoacidosis (the single most important risk factor), uncontrolled diabetes mellitus, neutropenia, haematological malignancy, organ/haematopoietic stem cell transplant, long-term corticosteroid or deferoxamine use, traumatic inoculation, burns, and malnutrition
    \item \textbf{Pathogenesis:} sporangiospores inhaled or inoculated $\to$ evade phagocytic clearance (impaired macrophages and neutrophils in the immunocompromised host) $\to$ germinate into hyphae $\to$ angioinvasive growth with thrombosis of blood vessels
    \item \textbf{Tissue destruction:} hyphae invade arterial walls, causing ischaemia, necrosis, and haemorrhage; this angioinvasion explains the characteristic black necrotic eschar and rapid local spread
    \item \textbf{Host iron/glucose interactions:} elevated free iron (deferoxamine, ketoacidosis) and hyperglycaemia promote fungal growth (Rhizopus uses a ferricrocin siderophore)
\end{itemize}

\section*{Clinical Features}
\begin{itemize}
    \item \textbf{Rhino-orbital-cerebral (most common, ~50\%):} facial pain, nasal congestion or black necrotic eschar on the palate/nasal mucosa, fever, headache, periorbital swelling, proptosis, ophthalmoplegia, facial anaesthesia, and vision loss; may progress to brain invasion
    \item \textbf{Pulmonary:} fever, cough, pleuritic chest pain, haemoptysis, dyspnoea; imaging may show consolidation, nodules, wedge-shaped infarcts, or cavitation
    \item \textbf{Cutaneous:} necrotic ulcer or eschar following trauma/burns, rapidly enlarging, painful, with surrounding erythema
    \item \textbf{Gastrointestinal:} abdominal pain, nausea, vomiting, haematemesis (commonly in premature neonates)
    \item \textbf{Disseminated:} haematogenous spread to brain, lungs, heart, kidneys; hypotension, multiorgan failure
    \item \textbf{Nonspecific:} fever not responding to antibiotics, malaise
\end{itemize}

\section*{Differential Diagnosis}
\begin{itemize}
    \item Other angioinvasive moulds: \textit{Aspergillus} (aspergillosis), \textit{Fusarium}, \textit{Scedosporium}, \textit{Lomentospora}
    \item Bacterial orbital cellulitis or cavernous sinus thrombosis (especially in diabetics)
    \item Sinusitis or nasal polyposis complicated by infection
    \item Malignancy (sinonasal tumour, lymphoma) causing facial swelling and necrosis
    \item Tuberculosis or other granulomatous disease in endemic areas
    \item Aseptic intracranial thrombophlebitis or stroke
\end{itemize}

\section*{When to Seek Medical Care}
Mucormycosis is a medical emergency requiring urgent specialist assessment. Patients with diabetes or immunosuppression who develop facial pain, nasal congestion with black necrotic tissue, orbital swelling, or cough with haemoptysis should seek care immediately.

\textbf{Call 000 or attend an emergency department for:}
\begin{itemize}
    \item Rapidly spreading facial or orbital swelling with visual disturbance or proptosis
    \item Black necrotic eschar in the nose, palate, or skin
    \item Severe headache with fever and neurological signs (altered consciousness, seizures)
    \item Haemoptysis or severe respiratory distress
    \item Signs of sepsis or shock
\end{itemize}

\section*{Diagnosis and Investigations}
\begin{itemize}
    \item \textbf{Histopathology:} biopsy of necrotic tissue demonstrating broad, non-septate (ribbon-like), right-angle branching hyphae within necrotic tissue and blood vessels (diagnostic)
    \item \textbf{Microscopy/culture:} direct microscopy with calcofluor white or KOH; culture on Sabouraud agar (may be negative despite disease)
    \item \textbf{Molecular:} PCR on fresh tissue for rapid species identification; increasingly used where available
    \item \textbf{Imaging:} CT or MRI of the sinuses/orbit/brain showing mucosal thickening, sinus opacification, bone erosion, and intracranial extension; chest CT for pulmonary disease
    \item \textbf{Endoscopic examination:} ENT nasal endoscopy with biopsy of necrotic mucosa
    \item \textbf{History/examination:} identify risk factors (diabetes, malignancy, transplant, immunosuppression, trauma)
\end{itemize}

\section*{Management}
\begin{itemize}
    \item \textbf{Emergency surgery:} urgent and aggressive surgical debridement of all necrotic tissue (rhino-orbital-cerebral involvement often requires orbital exenteration if the orbit is invaded); repeat debridement frequently
    \item \textbf{Antifungal therapy:} first-line liposomal amphotericin B (e.g., 5--10 mg/kg/day IV) commenced empirically and urgently; high-dose posaconazole or isavuconazole as alternatives or for step-down once stable
    \item \textbf{Reversal of predisposing factors:} strict glycaemic control, correction of ketoacidosis, reduce or stop corticosteroids and deferoxamine, restore neutrophil counts (G-CSF where appropriate)
    \item \textbf{Adjuncts:} surgery plus antifungals is standard; hyperbaric oxygen and cytokine therapy are of unproven benefit
    \item \textbf{Mild/limited cutaneous disease:} surgical excision plus oral antifungal (posaconazole/isavuconazole) may suffice
    \item \textbf{Supportive care:} ICU, respiratory and circulatory support for severe/disseminated disease
\end{itemize}

\section*{Complications}
\begin{itemize}
    \item Cerebral invasion with brain abscess, stroke, or meningitis
    \item Loss of vision or an eye (orbital exenteration)
    \item Cavernous sinus thrombosis and internal carotid artery invasion/aneurysm
    \item Massive haemoptysis from pulmonary angioinvasion
    \item Sepsis, septic shock, and multiorgan failure
    \item Disfigurement needing reconstructive surgery
\end{itemize}

\section*{Prognosis and Outcomes}
Mucormycosis carries a grave prognosis, especially when diagnosis is delayed. Mortality is approximately 20--30\% for localised cutaneous disease, 30--50\% for rhino-orbital-cerebral disease, and 50--90\% for pulmonary and disseminated disease. Survival improves substantially with early diagnosis, urgent surgical debridement, appropriate antifungal therapy, and reversal of the underlying predisposing condition (especially diabetic ketoacidosis). Even survivors may have significant long-term morbidity (vision loss, facial disfigurement, neurological deficit).

\section*{Prevention and Self-Care}
\begin{itemize}
    \item Strict glycaemic control in patients with diabetes, especially during illness or steroid use
    \item Judicious use of corticosteroids and prompt tapering once clinically indicated
    \item Neutropenic precautions (protective isolation, prophylaxis where appropriate) in high-risk transplant/chemotherapy patients
    \item Avoidance of environmental mould exposure in at-risk patients (dust, soil, construction sites, decaying vegetation)
    \item Prompt and thorough wound care and debridement of contaminated traumatic wounds in immunocompromised patients
    \item Early ENT assessment of sinus symptoms in immunocompromised patients
\end{itemize}

\section*{Sources and Further Reading}
\begin{itemize}
    \item Healthdirect Australia. \textit{Mucormycosis (black fungus).} https://www.healthdirect.gov.au/
    \item Centers for Disease Control and Prevention (CDC). \textit{Mucormycosis.} https://www.cdc.gov/fungal/diseases/mucormycosis/
    \item Therapeutic Guidelines (Australia). \textit{Antifungal guidelines.} https://www.tg.org.au/
    \item National Health and Medical Research Council (NHMRC). \textit{Invasive fungal infection guidance} (via relevant societies)
    \item Cornely OA et al. \textit{Global guideline for the diagnosis and management of mucormycosis.} Lancet Infect Dis 2019.
    \item Clinical Excellence Commission NSW. \textit{Antimicrobial stewardship resources.}
\end{itemize}

\clearpage
